% Arquitectura Multi-Agente con IA Generativa: Plasticidad de la Estructura
% ArXiv Submission - December 2025
\documentclass[11pt,a4paper]{article}

% ============================================================================
% PACKAGES
% ============================================================================

% Encoding and fonts
\usepackage[utf8]{inputenc}
\usepackage[T1]{fontenc}
\usepackage{lmodern}

% Mathematics
\usepackage{amsmath,amssymb,amsthm}
\usepackage{mathtools}

% Graphics and figures
\usepackage{graphicx}
\usepackage{float}
\usepackage{subcaption}

% Fix transparent package compatibility (must be before svg)
\makeatletter
\providecommand{\IfPDFManagementActiveTF}[2]{#2}
\providecommand{\IfPDFManagementActiveF}[1]{#1}
\makeatother

\usepackage[inkscapelatex=false,inkscapeexe={"C:/Program Files/Inkscape/bin/inkscape.exe"}]{svg}
\svgpath{{../diagrams/}}
\graphicspath{{figures/}}

% Tables
\usepackage{booktabs}
\usepackage{multirow}
\usepackage{array}
\usepackage{longtable}

% Colors and hyperlinks
\usepackage{xcolor}
\usepackage[colorlinks=true,linkcolor=blue,citecolor=blue,urlcolor=blue]{hyperref}

% Code listings
\usepackage{listings}
\usepackage{algorithm}
\usepackage{algpseudocode}

% Layout
\usepackage[margin=1in]{geometry}
\usepackage{setspace}
\onehalfspacing

% Bibliography
\usepackage[numbers,sort&compress]{natbib}

% Misc
\usepackage{enumitem}
\usepackage{xspace}

% ============================================================================
% THEOREM ENVIRONMENTS
% ============================================================================

\newtheorem{theorem}{Theorem}[section]
\newtheorem{lemma}[theorem]{Lemma}
\newtheorem{proposition}[theorem]{Proposition}
\newtheorem{corollary}[theorem]{Corollary}
\theoremstyle{definition}
\newtheorem{definition}{Definition}[section]
\newtheorem{example}{Example}[section]
\theoremstyle{remark}
\newtheorem{remark}{Remark}[section]

% ============================================================================
% CUSTOM COMMANDS
% ============================================================================

\newcommand{\lumen}{\textsc{Lumen}\xspace}
\newcommand{\pema}{\textsc{PEMA}\xspace}
\newcommand{\sgrag}{\textsc{SG-RAG}\xspace}

% Code style
\definecolor{codebg}{RGB}{248,248,248}
\definecolor{codegreen}{RGB}{0,128,0}
\definecolor{codegray}{RGB}{128,128,128}
\definecolor{codepurple}{RGB}{128,0,128}

\lstdefinestyle{pythonstyle}{
    backgroundcolor=\color{codebg},
    basicstyle=\ttfamily\small,
    breaklines=true,
    captionpos=b,
    commentstyle=\color{codegreen},
    keywordstyle=\color{codepurple}\bfseries,
    stringstyle=\color{codegray},
    showstringspaces=false,
    frame=single,
    framerule=0.5pt,
    numbers=left,
    numberstyle=\tiny\color{codegray},
    tabsize=4
}
\lstset{style=pythonstyle}

% ============================================================================
% DOCUMENT INFO
% ============================================================================

\title{\textbf{Multi-Agent Architecture with Generative AI:\\Structural Plasticity}\\[0.5em]
\large Design Patterns and Formal Foundations\\for Enterprise Intelligent Systems}

\author{
    Ariel Edgardo Levy\\
    \texttt{ariel.e.levy@gmail.com}\\
    Posadas, Argentina
}

\date{December 2025}

% ============================================================================
% DOCUMENT
% ============================================================================

\begin{document}

\maketitle

% ----------------------------------------------------------------------------
% ABSTRACT
% ----------------------------------------------------------------------------

\begin{abstract}
This paper presents \lumen, a multi-agent agentic AI platform designed for Business Intelligence environments. We introduce three main contributions: (1) \textbf{Functional Agency}, a formal framework defining the necessary and sufficient conditions for agentic behavior in production systems; (2) \textbf{\pema} (Plastic Evolving Multi-Agent), a novel architecture that applies Hebbian learning principles for dynamic inter-agent routing with formal convergence guarantees; and (3) a catalog of \textbf{12 reusable design patterns} for constructing robust enterprise AI systems.

The platform integrates semantic routing, multi-layer RAG with gap detection (\sgrag), and anti-hallucination patterns, achieving 94\% task success rate in production environments while maintaining sub-2-second latency for 85\% of queries. Our empirical evaluation across three enterprise deployments demonstrates significant improvements in both efficiency (40\% cost reduction via hybrid routing) and reliability (hallucination reduction from 12\% to 0.3\%).

Beyond the technical implementation, we formalize the theoretical foundations of multi-agent coordination, including Theorem 10.1 establishing variance bounds under plasticity, and propose PEMA-Bench, a standardized evaluation protocol for adaptive multi-agent systems. All architectural patterns and core implementations are open-source, enabling reproducibility and community extension.

\vspace{0.5em}
\noindent\textbf{Keywords:} Multi-agent systems, Agentic AI, Business Intelligence, RAG, Design Patterns, Hebbian Learning, LLM Orchestration
\end{abstract}

\newpage

% ============================================================================
% PART I: FOUNDATIONS
% ============================================================================

\part{Theoretical Foundations}

% ----------------------------------------------------------------------------
% CHAPTER 1: INTRODUCTION
% ----------------------------------------------------------------------------

\section{Introduction}
\label{sec:introduction}

\subsection{The Emergence of Agentic AI}

The evolution of artificial intelligence systems has followed a clear progression: from passive models that responded to isolated queries~\cite{brown2020language}, to reactive systems capable of contextual responses~\cite{openai2023gpt4}, to the current frontier of \textit{agentic systems}---autonomous entities capable of planning, executing, and learning from multi-step interactions with their environment~\cite{wooldridge1995intelligent}.

\begin{figure}[H]
    \centering
    \includesvg[width=0.9\textwidth]{1.1_evolucion_sistemas_ia_en}
    \caption{Evolution of AI systems: from passive to agentic}
    \label{fig:evolution-ai}
\end{figure}

This transition is not merely incremental; it represents a fundamental change in how we conceive the relationship between AI and the tasks it performs. While traditional systems act as sophisticated tools that amplify human capabilities, agentic systems begin to exhibit characteristics previously reserved for autonomous entities: the ability to set goals, plan strategies, and adapt their behavior based on accumulated experience.

In the Business Intelligence (BI) context, this transition has particular relevance. BI environments are characterized by:

\begin{itemize}
    \item \textbf{Semantic complexity}: Multiple domains (financial, operational, commercial) with specific vocabularies
    \item \textbf{Structural heterogeneity}: Diverse data sources with incompatible schemas
    \item \textbf{Temporal dynamism}: Models and metrics that evolve with business needs
    \item \textbf{High stakes}: Decisions based on incorrect analysis can have significant economic impact
\end{itemize}

\subsection{Towards a Formal Definition of Agency}

The term ``agent'' has been used loosely in the AI community, often referring to any system that performs actions~\cite{wooldridge1995intelligent}. However, for rigorous system design, we need a more precise definition that distinguishes truly agentic systems from sophisticated pipelines.

\begin{figure}[H]
    \centering
    \includesvg[width=0.8\textwidth]{1.2_agencia_funcional_en}
    \caption{The three dimensions of Functional Agency}
    \label{fig:functional-agency}
\end{figure}

\begin{definition}[Functional Agency]
\label{def:functional-agency}
A system $\mathcal{A}$ exhibits \textbf{functional agency} if and only if it satisfies three conditions:

\begin{enumerate}
    \item \textbf{Perception-Action Loop}: $\mathcal{A}$ maintains a continuous cycle $\mathcal{P}: \mathcal{E} \rightarrow \mathcal{S}$ and $\mathcal{A}: \mathcal{S} \rightarrow \mathcal{E}$ where $\mathcal{E}$ is the environment and $\mathcal{S}$ is the internal state space.

    \item \textbf{Goal-Directed Behavior}: There exists an objective function $\mathcal{O}: \mathcal{S} \times \mathcal{E} \rightarrow \mathbb{R}$ that $\mathcal{A}$ actively seeks to maximize.

    \item \textbf{Adaptive Autonomy}: $\mathcal{A}$ can modify its strategy $\pi: \mathcal{S} \rightarrow \mathcal{A}$ based on accumulated experience without explicit reprogramming.
\end{enumerate}
\end{definition}

This definition excludes static pipelines (failing condition 3) and pure LLM systems without memory (failing condition 1), while including systems that genuinely learn and adapt from their interactions.

\subsection{The Agentic Cycle}

Based on Definition~\ref{def:functional-agency}, we can formalize the fundamental operational cycle of an agentic system:

\begin{figure}[H]
    \centering
    \includesvg[width=0.85\textwidth]{1.3_ciclo_agentico_en}
    \caption{The agentic cycle: perceive, reason, act, learn}
    \label{fig:agentic-cycle}
\end{figure}

\begin{definition}[Agentic Cycle]
\label{def:agentic-cycle}
The \textbf{agentic cycle} is defined as the tuple $\mathcal{C} = (\mathcal{P}, \mathcal{R}, \mathcal{X}, \mathcal{L})$ where:

\begin{itemize}
    \item $\mathcal{P}$: \textit{Perceive} --- Capture and process environmental signals
    \item $\mathcal{R}$: \textit{Reason} --- Deliberate using internal models and accumulated knowledge
    \item $\mathcal{X}$: \textit{Execute} --- Perform actions that modify the environment
    \item $\mathcal{L}$: \textit{Learn} --- Update internal models based on observed results
\end{itemize}

The cycle satisfies continuity: $\mathcal{L}_t \rightarrow \mathcal{P}_{t+1}$, where each learning iteration informs the next perception phase.
\end{definition}

\subsection{Hypothesis and Contributions}

This paper develops the following central hypothesis:

\begin{quote}
\textit{Multi-agent systems with formal agency guarantees, specialized routing, and plastic adaptation mechanisms can achieve superior performance in complex BI environments compared to monolithic approaches or static pipelines.}
\end{quote}

To test this hypothesis, we present:

\begin{enumerate}
    \item \textbf{Functional Agency Framework}: Formal definitions and necessary conditions for production agentic systems (Definition~\ref{def:functional-agency})

    \item \textbf{\pema Architecture}: Hebbian learning for inter-agent routing with convergence guarantees (Section~\ref{sec:pema})

    \item \textbf{Design Pattern Catalog}: 12 reusable patterns for enterprise AI (throughout the document)

    \item \textbf{Complete Implementation}: \lumen, a production-ready platform (Sections~\ref{sec:architecture}--\ref{sec:observability})

    \item \textbf{Empirical Evaluation}: Results from three enterprise deployments (Section~\ref{sec:evaluation})
\end{enumerate}

\subsection{Document Structure}

The document is organized in three parts:

\textbf{Part I: Foundations} (Sections 1--2) establishes the theoretical foundations, including the formal definition of agency, taxonomy of multi-agent architectures, and analysis of when multi-agent systems are appropriate.

\textbf{Part II: Architecture} (Sections 3--4 and Capabilities 1--10) details the complete \lumen implementation, from routing patterns to the \pema learning system.

\textbf{Part III: Evaluation} (Sections 5--6) presents empirical results, discussion of limitations, and future directions.

\subsection{Current Challenges in Agentic AI}

Despite advances in language models, building reliable agentic systems faces fundamental challenges:

\begin{figure}[H]
    \centering
    \vspace{0.5em}
    \includesvg[width=0.95\textwidth]{1.6_desafios_actuales_en}
    \vspace{0.5em}
    \caption{Current challenges in building agentic AI systems}
    \label{fig:current-challenges}
\end{figure}

\begin{enumerate}
    \item \textbf{The Reliability Problem}: LLMs are probabilistic by nature~\cite{ji2023survey}; ensuring consistent behavior in critical applications remains an open challenge.

    \item \textbf{The Coordination Problem}: How to orchestrate multiple specialized agents without creating emergent chaotic behavior?

    \item \textbf{The Learning Problem}: How to enable continuous adaptation without catastrophic forgetting~\cite{kirkpatrick2017overcoming} or dangerous drift?

    \item \textbf{The Observability Problem}: How to understand and debug systems whose behavior emerges from multiple interacting agents?
\end{enumerate}

This paper addresses these challenges through a combination of formal foundations, robust patterns, and extensive empirical validation.

\subsection{Horizons of Agentic AI}

Looking to the future, we identify three main development horizons:

\begin{figure}[H]
    \centering
    \includesvg[width=0.9\textwidth]{1.7_horizonte_ia_agentica_en}
    \caption{Development horizons of agentic AI}
    \label{fig:horizons}
\end{figure}

\begin{description}
    \item[Short-term (1--2 years):] Consolidation of design patterns, standardization of evaluation frameworks, expansion into regulated vertical domains.

    \item[Medium-term (3--5 years):] Emergence of meta-agents capable of composing other agents, formal verification of behavioral guarantees, integration with symbolic systems for hybrid reasoning.

    \item[Long-term (5+ years):] Truly autonomous systems with minimal supervision, continuous learning in open environments, self-repairing and self-evolving architectures.
\end{description}

% ----------------------------------------------------------------------------
% CHAPTER 2: MULTI-AGENT ARCHITECTURES
% ----------------------------------------------------------------------------

\section{Multi-Agent Architectures}
\label{sec:multiagent}

\subsection{Agent Architecture Taxonomy}

The design space for multi-agent systems can be organized along two fundamental axes: \textit{coupling} (how tightly agents are connected) and \textit{control} (how decisions are coordinated).

\subsubsection{Hierarchical Architecture}

\begin{figure}[H]
    \centering
    \includesvg[width=0.7\textwidth]{2.2.1_arquitectura_jerarquica_en}
    \caption{Hierarchical architecture with supervisor coordination}
    \label{fig:hierarchical}
\end{figure}

In hierarchical architectures, a supervisor agent coordinates specialized workers. This pattern is ideal when:
\begin{itemize}
    \item There is a clear task decomposition
    \item Global visibility is required for coordination
    \item A single point of control is acceptable
\end{itemize}

\subsubsection{Swarm Architecture}

\begin{figure}[H]
    \centering
    \includesvg[width=0.7\textwidth]{2.2.2_arquitectura_swarm_en}
    \caption{Swarm architecture with emergent coordination}
    \label{fig:swarm}
\end{figure}

Swarm systems exhibit coordination through local rules without central control. Advantages include:
\begin{itemize}
    \item High resilience to individual failures
    \item Natural scalability
    \item Emergent behavior capable of solving complex problems
\end{itemize}

\subsubsection{Meta-Learning Architecture}

\begin{figure}[H]
    \centering
    \includesvg[width=0.7\textwidth]{2.2.3_arquitectura_meta_aprendizaje_en}
    \caption{Meta-learning architecture with routing adaptation}
    \label{fig:meta-learning}
\end{figure}

Meta-learning systems include a meta-agent that observes agent performance and adjusts routing strategies~\cite{wu2023autogen,hong2023metagpt}. This is the basis for our \pema architecture.

\subsubsection{Additional Architectures}

For completeness, we describe additional patterns:

\paragraph{Modular Architecture} Interchangeable specialized modules with standardized interfaces.

\paragraph{Evolutionary Architecture} Agent population with selection and recombination.

\paragraph{Orchestrator-Worker} Explicit orchestration with task dispatch.

\paragraph{Generator-Critic} Adversarial pattern for quality assurance.

\paragraph{Blackboard Architecture} Shared memory for collaborative problem-solving.

\paragraph{Sequential Architecture} Linear pipelines with handoffs.

\paragraph{Parallel Architecture} Concurrent execution with result aggregation.

\paragraph{Reflexive Architecture} Self-observation and self-modification capabilities.

\subsection{When to Use Multi-Agent Systems}

Not every problem requires multi-agent complexity. We provide a decision framework:

\begin{figure}[H]
    \centering
    \includesvg[width=0.8\textwidth]{2.4.3_single_vs_multiagent_en}
    \caption{Decision framework: single-agent vs multi-agent}
    \label{fig:single-vs-multi}
\end{figure}

\textbf{Use single-agent when:}
\begin{itemize}
    \item The domain is homogeneous
    \item Task complexity is manageable
    \item Simplicity is paramount
\end{itemize}

\textbf{Use multi-agent when:}
\begin{itemize}
    \item Multiple specialized domains exist
    \item Tasks require different capabilities
    \item Scalability and resilience are critical
\end{itemize}

% ----------------------------------------------------------------------------
% CHAPTER 3: RELATED WORK
% ----------------------------------------------------------------------------

\section{Related Work}
\label{sec:related-work}

\subsection{Orchestration Frameworks}

Several frameworks have emerged for multi-agent LLM systems:

\textbf{LangChain/LangGraph}~\cite{langchain2024}: Pioneer framework introducing ``chains'' for LLM composition. LangGraph extends this with state graphs for complex workflows.

\textbf{Microsoft AutoGen}~\cite{wu2023autogen}: Multi-agent conversation framework where agents interact via natural dialogue.

\textbf{CrewAI}: Simplified framework with ``crew'' metaphor of agents with defined roles.

\textbf{OpenAI Swarm}: Minimalist experimental framework for agent ``handoffs''.

\subsection{Conversational BI Systems}

Existing BI solutions with natural language interfaces include Power BI Q\&A, ThoughtSpot, and Tableau Ask Data. However, these systems are limited to simple queries and lack agentic architecture.

\subsection{Adaptive Multi-Agent Systems}

A critical gap in previous work is \textbf{structural adaptation}. Existing approaches present static topologies:

\textbf{DyLAN}~\cite{liu2024dylan}: Proposes dynamic agent networks but operates at agent selection per task, not cumulative trust weights.

\textbf{MetaGPT}~\cite{hong2023metagpt}: Multi-agent framework with well-defined roles but fixed topology that doesn't learn from experiences.

\textbf{Difference with MARL}: Multi-Agent Reinforcement Learning operates on policies learned end-to-end via gradients, requiring millions of training interactions. \pema operates on pre-trained agents (LLMs), adjusting only \textit{inter-agent trust weights} via local Hebbian rules---allowing adaptation with orders of magnitude less data.

\subsection{Comparative Positioning}

Table~\ref{tab:comparison} summarizes how \lumen addresses gaps in existing approaches:

\begin{table}[H]
\centering
\small
\begin{tabular}{lccccc}
\toprule
\textbf{System} & \textbf{Multi-Agent} & \textbf{BI Native} & \textbf{Adaptive} & \textbf{Memory} & \textbf{Anti-Halluc.} \\
\midrule
LangGraph & \checkmark & -- & -- & Partial & -- \\
AutoGen & \checkmark & -- & -- & Partial & -- \\
CrewAI & \checkmark & -- & -- & -- & -- \\
Power BI Q\&A & -- & \checkmark & -- & -- & -- \\
ThoughtSpot & -- & \checkmark & -- & -- & -- \\
\textbf{\lumen} & \checkmark & \checkmark & \checkmark & \checkmark & \checkmark \\
\bottomrule
\end{tabular}
\caption{Comparative positioning of \lumen vs. related systems}
\label{tab:comparison}
\end{table}

\subsection{Identified Gaps}

Analysis of previous works reveals gaps this work addresses:

\begin{enumerate}
    \item \textbf{BI Integration Gap}: Agent frameworks aren't optimized for BI; BI solutions don't have agentic architecture.
    \item \textbf{Memory Gap}: No system combines structured persistent memory between specialized agents.
    \item \textbf{Routing Gap}: Approaches are either too simple (keywords) or too expensive (always LLM).
    \item \textbf{Adaptation Gap}: None addresses the \textit{adaptation barrier}---inability to modify coordination patterns based on observed results.
\end{enumerate}

% ============================================================================
% PART II: ARCHITECTURE
% ============================================================================

\part{Architecture and Implementation}

% ----------------------------------------------------------------------------
% CHAPTER 3: LUMEN ARCHITECTURE
% ----------------------------------------------------------------------------

\section{\lumen Architecture Overview}
\label{sec:architecture}

\lumen is structured as a layered system with clear separation of concerns, following Clean Architecture principles~\cite{martin2017clean,fowler2002patterns}:

\begin{figure}[H]
    \centering
    \includesvg[width=0.95\textwidth]{P2_lumen_arquitectura_en}
    \caption{Complete \lumen architecture showing all layers and components}
    \label{fig:lumen-architecture}
\end{figure}

\subsection{Architectural Layers}

\begin{enumerate}
    \item \textbf{Interface Layer}: HTTP/WebSocket APIs, authentication, session management
    \item \textbf{Routing Layer}: Two-layer routing (keyword + semantic), intent classification
    \item \textbf{Agent Layer}: Specialized agents with tool access and memory
    \item \textbf{Knowledge Layer}: RAG, vector stores, semantic search
    \item \textbf{Integration Layer}: External systems, databases, BI platforms
    \item \textbf{Infrastructure Layer}: Observability, caching, persistence
\end{enumerate}

% ----------------------------------------------------------------------------
% CAPABILITIES 1-10
% ----------------------------------------------------------------------------

\section{Capability 1: Memory and Persistence}
\label{sec:memory}

\subsection{Memory Architecture}

Effective agentic systems require sophisticated memory management across multiple time horizons and abstraction levels.

\begin{figure}[H]
    \centering
    \includesvg[width=0.85\textwidth]{C1_gestion_memoria_en}
    \caption{Multi-tier memory architecture}
    \label{fig:memory}
\end{figure}

\begin{definition}[Memory Tiers]
\lumen implements three memory tiers:
\begin{itemize}
    \item \textbf{Working Memory}: In-session context (Redis, TTL: session duration)
    \item \textbf{Episodic Memory}: Cross-session user history (SQL, TTL: configurable)
    \item \textbf{Semantic Memory}: Domain knowledge (Vector store, persistent)
\end{itemize}
\end{definition}

\section{Capability 2: Semantic Routing}
\label{sec:routing}

\subsection{Two-Layer Routing Pattern}

A key innovation is our \textbf{Two-Layer Routing} pattern that balances cost and accuracy:

\begin{figure}[H]
    \centering
    \includesvg[width=0.85\textwidth]{C2_semantic_routing_en}
    \caption{Two-layer routing: keywords first, then semantic}
    \label{fig:routing}
\end{figure}

This pattern achieves:
\begin{itemize}
    \item \textbf{5x cost reduction} vs pure LLM routing
    \item \textbf{Sub-100ms} latency for 85\% of queries
    \item \textbf{98\% accuracy} on intent classification
\end{itemize}

\section{Capability 3: Hybrid RAG}
\label{sec:rag}

\subsection{SG-RAG: Speculative Gap RAG}

Traditional RAG systems~\cite{lewis2020retrieval,gao2023retrieval} fail silently when information is missing. \sgrag addresses this:

\begin{figure}[H]
    \centering
    \includesvg[width=0.85\textwidth]{C3_rag_lumen_en}
    \caption{\sgrag architecture with gap detection}
    \label{fig:sgrag}
\end{figure}

\begin{definition}[Speculative Gap RAG]
\sgrag extends traditional RAG with three mechanisms:
\begin{enumerate}
    \item \textbf{Confidence Scoring}: Each retrieval includes confidence metrics
    \item \textbf{Gap Detection}: Low-confidence retrievals are flagged as potential gaps
    \item \textbf{Gap Persistence}: Detected gaps are stored for documentation improvement
\end{enumerate}
\end{definition}

\subsubsection{Gap Detection Algorithm}

The core innovation of \sgrag lies in its \textbf{speculative gap detection}, which identifies when the knowledge base lacks sufficient information to answer a query. A gap is detected when:

\begin{equation}
\text{Gap}(q, D) = \mathbb{1}\left[\max_{d \in D} \text{sim}(q, d) < \tau_{\text{conf}} \lor \text{coverage}(q, D) < \tau_{\text{cov}}\right]
\end{equation}

where $q$ is the query, $D$ is the document corpus, $\tau_{\text{conf}} = 0.75$ is the confidence threshold, and $\tau_{\text{cov}}$ is the coverage threshold.

\textbf{Gap detection criteria:}
\begin{enumerate}
    \item \textbf{Low similarity score}: Best matching document has cosine similarity $< 0.75$
    \item \textbf{High variance}: Top-$k$ results have high score variance ($\sigma^2 > 0.1$), indicating ambiguous matches
    \item \textbf{Semantic mismatch}: Query entities not present in retrieved context (entity coverage $< 0.5$)
    \item \textbf{Answer uncertainty}: LLM response contains hedging phrases (``I don't have information'', ``not specified'')
\end{enumerate}

When a gap is detected, \sgrag:
\begin{itemize}
    \item Returns a transparent response acknowledging the limitation
    \item Logs the gap with query, timestamp, and attempted retrievals
    \item Triggers asynchronous notification to documentation maintainers
    \item Caches the gap pattern for future detection optimization
\end{itemize}

This ``fail loudly'' approach contrasts with traditional RAG systems that hallucinate when information is missing, achieving a 96\% reduction in false-positive answers in our production deployment.

\section{Capability 4: Agent Orchestration}
\label{sec:orchestration}

\subsection{Agent-as-Tool Pattern}

The Agent-as-Tool pattern enables seamless composition of specialized agents, extending the tool-use paradigm~\cite{schick2024toolformer,patil2023gorilla}:

\begin{figure}[H]
    \centering
    \includesvg[width=0.8\textwidth]{C4_agent_as_tool_en}
    \caption{Agent-as-Tool pattern for composable agents}
    \label{fig:agent-as-tool}
\end{figure}

\subsection{Handoff Pattern}

The Handoff pattern allows an agent to \textbf{completely} transfer control to another agent when it detects that the task exceeds its expertise or requires specialized capabilities.

\begin{definition}[Handoff]
A handoff $H$ from agent $A_i$ to agent $A_j$ is a tuple $H = (\text{trigger}, \text{target}, \text{context}, \text{query})$ where:
\begin{itemize}
    \item \textbf{trigger}: Condition that triggers the transfer (e.g., ``requires DAX execution'')
    \item \textbf{target}: Target agent with required expertise
    \item \textbf{context}: State and metadata transferred (session\_id, thread, history)
    \item \textbf{query}: Original user input preserved
\end{itemize}
\end{definition}

\begin{figure}[H]
    \centering
    \includesvg[width=0.8\textwidth]{C4_handoff_pattern_en}
    \caption{Handoff Pattern---Complete Transfer of Control}
    \label{fig:handoff-pattern}
\end{figure}

\subsubsection{Types of Handoff}

\begin{table}[H]
\centering
\begin{tabular}{lll}
\toprule
\textbf{Type} & \textbf{Description} & \textbf{Example} \\
\midrule
Expertise-based & Agent lacks specialized knowledge & GeneralAgent $\rightarrow$ DAXAgent \\
Capability-based & Required tools not available & PowerBIAgent $\rightarrow$ FabricAgent \\
Context-based & Context indicates another domain & DocumentAgent $\rightarrow$ SearchAgent \\
Planning-based & Multi-step query needs orchestration & FabricAgent $\rightarrow$ PlanningAgent \\
\bottomrule
\end{tabular}
\caption{Types of handoff in \lumen}
\label{tab:handoff-types}
\end{table}

\subsubsection{Detection Mechanism}

Handoff is triggered through two complementary mechanisms:

\begin{enumerate}
    \item \textbf{Proactive detection by workflow}: The workflow evaluates if the current message requires an agent change based on domain keywords.

    \item \textbf{Explicit request by agent}: When an agent internally detects it cannot complete the task, it requests handoff to an agent that can orchestrate the complete sequence.
\end{enumerate}

\subsubsection{Context Preservation}

During handoff the following are preserved:
\begin{itemize}
    \item \textbf{Conversation thread}: Target agent continues in the same multi-turn thread
    \item \textbf{Session metadata}: user\_id, session\_id, available tokens
    \item \textbf{Original query}: User input unmodified
    \item \textbf{Handoff reason}: For logging and debugging
\end{itemize}

\subsection{Agent Hierarchy}

\begin{figure}[H]
    \centering
    \vspace{0.5em}
    \includesvg[width=0.95\textwidth]{C4_jerarquia_agentes_en}
    \vspace{0.5em}
    \caption{\lumen agent hierarchy}
    \label{fig:agent-hierarchy}
\end{figure}

\section{Capability 5: Data Source Integration}
\label{sec:integration}

\lumen integrates with multiple BI data sources through a unified abstraction layer:

\begin{itemize}
    \item \textbf{Power BI}: XMLA endpoints, semantic models
    \item \textbf{SQL Databases}: Direct query with schema introspection
    \item \textbf{REST APIs}: Generic connector framework
    \item \textbf{File Sources}: Excel, CSV, Parquet
\end{itemize}

\section{Capability 6: Visualization and Output}
\label{sec:visualization}

\subsection{Rendering Pipeline}

\begin{figure}[H]
    \centering
    \includesvg[width=0.85\textwidth]{C6_rendering_pipeline_en}
    \caption{Visualization rendering pipeline}
    \label{fig:rendering}
\end{figure}

The rendering pipeline supports multiple output formats:
\begin{itemize}
    \item Interactive charts (ECharts, Plotly)
    \item Static exports (PNG, SVG, PDF)
    \item Embedded Power BI visuals
    \item Markdown tables
\end{itemize}

\section{Capability 7: Security and Anti-Hallucination}
\label{sec:security}

\subsection{HallucinationGuard Pattern}

\begin{figure}[H]
    \centering
    \includesvg[width=0.9\textwidth]{C7_hallucination_guard_en}
    \caption{HallucinationGuard Pattern: validates agent outputs against grounded sources, reducing hallucinations from 12\% to 0.3\%}
    \label{fig:hallucination-guard}
\end{figure}

This pattern reduced hallucinations from 12\% to 0.3\% in production, addressing a key challenge identified in~\cite{huang2023survey,manakul2023selfcheckgpt}.

\section{Capability 8: Scalability}
\label{sec:scalability}

\subsection{Horizontal Scaling Architecture}

\begin{figure}[H]
    \centering
    \includesvg[width=\textwidth]{C8_arquitectura_aks_en}
    \caption{Kubernetes-based horizontal scaling}
    \label{fig:scaling}
\end{figure}

Key scalability features:
\begin{itemize}
    \item \textbf{Stateless agents}: All state externalized to Redis/SQL
    \item \textbf{Auto-scaling}: HPA based on queue depth and latency
    \item \textbf{Connection pooling}: Efficient resource utilization
\end{itemize}

\section{Capability 9: Observability}
\label{sec:observability}

\subsection{Observability Stack}

\begin{figure}[H]
    \centering
    \includesvg[width=0.85\textwidth]{C8_observability_stack_en}
    \caption{Complete observability stack}
    \label{fig:observability}
\end{figure}

\lumen implements \textbf{agent-aware observability}, extending traditional monitoring with agentic-specific telemetry:

\begin{enumerate}
    \item \textbf{Agent Decision Traces}: Complete reasoning chains from query to response, including tool invocations, handoffs, and intermediate states---enabling root-cause analysis of agent behavior

    \item \textbf{Semantic Metrics}: Beyond latency/throughput, we track task success rate, hallucination incidents, gap detections, and routing accuracy per agent type

    \item \textbf{Token Economics}: Real-time tracking of token consumption per agent, model tier usage (GPT-5.2 vs GPT-5.2 nano), and cost attribution by conversation thread

    \item \textbf{Correlation Context}: Every request carries a \texttt{conversation\_id} propagated across all agents, tools, and async workers---enabling end-to-end trace reconstruction
\end{enumerate}

Implementation combines two complementary approaches:
\begin{itemize}
    \item \textbf{Real-time monitoring}: Redis Streams for live log aggregation, enabling operators to tail agent execution across distributed workers through a custom viewer (\texttt{logmon})
    \item \textbf{Historical analysis}: All interactions persisted to SQL database with full conversation context, enabling behavioral analysis, pattern detection, and continuous improvement of routing decisions
\end{itemize}

\section{Capability 10: PEMA --- Plastic Evolving Multi-Agent}
\label{sec:pema}

\subsection{Motivation}

Static routing rules cannot adapt to changing usage patterns. \pema addresses this through Hebbian learning principles~\cite{hebb1949organization} applied to agent coordination.

\begin{figure}[H]
    \centering
    \includesvg[width=0.85\textwidth]{10_pema_evolucion_en}
    \caption{Evolution from static to plastic systems in \pema}
    \label{fig:pema-evolution}
\end{figure}

\begin{figure}[H]
    \centering
    \includesvg[width=0.85\textwidth]{10_pema_trust_graph_en}
    \caption{Trust graph structure in \pema}
    \label{fig:trust-graph}
\end{figure}

\subsection{Hebbian Learning for Agent Routing}

\begin{definition}[Hebbian Weight Update]
For intent type $i$ and agent $j$, the trust weight $w_{ij}$ is updated:
\begin{equation}
w_{ij}(t+1) = w_{ij}(t) + \eta \cdot \text{success}_{ij}(t) - \lambda \cdot w_{ij}(t)
\end{equation}
where $\eta$ is the learning rate and $\lambda$ is the decay factor.
\end{definition}

\begin{remark}[Intuition: Hebbian Learning]
This rule embodies the principle \emph{``neurons that fire together, wire together.''} When agent $j$ successfully handles intent $i$ (success $> 0$), their connection strengthens. When it fails (success $< 0$), trust decreases. The decay term $\lambda w_{ij}$ prevents unbounded growth and allows the system to ``forget'' outdated patterns, enabling adaptation to changing agent capabilities.
\end{remark}

\subsection{Formal Guarantees}

\begin{theorem}[Variance Bound with Plasticity]
\label{thm:variance-bound}
Under \pema with learning rate $\eta > 0$ and decay $\lambda \in (0,1)$, the variance of routing decisions is bounded:
\begin{equation}
\text{Var}[\pi_t] \leq \frac{\eta^2 \sigma^2}{\lambda(2-\lambda)}
\end{equation}
where $\sigma^2$ is the variance of success signals.
\end{theorem}

\begin{proof}
See Appendix~\ref{sec:proofs}.
\end{proof}

\begin{remark}[Intuition]
The bound reveals a \textbf{stability-adaptability trade-off}: smaller $\eta$ reduces variance but slows learning; larger $\lambda$ increases forgetting but improves stability. The denominator $\lambda(2-\lambda)$ approaches its maximum of 1 as $\lambda \to 1$, so variance is minimized with aggressive decay.
\end{remark}

\begin{remark}[Assumptions]
Theorem~\ref{thm:variance-bound} requires: (i) success signals $s_t$ are i.i.d.\ with bounded variance $\sigma^2$; (ii) the intent distribution is stationary over the adaptation window; (iii) agent capabilities remain constant during learning. In practice, we use rolling windows of 100 interactions to approximate stationarity.
\end{remark}

\begin{corollary}[Practical Bound]
For typical hyperparameters $\eta = 0.01$, $\lambda = 0.1$, and $\sigma^2 \approx 0.25$:
\[
\text{Var}[\pi] \leq \frac{(0.01)^2 \times 0.25}{0.1 \times 1.9} = \frac{0.000025}{0.19} \approx 0.00013
\]
This ensures routing fluctuations below $1.2\%$, providing highly stable agent selection.
\end{corollary}

\subsection{PEMA-Bench: Evaluation Protocol}

We propose PEMA-Bench with 9 metrics across 3 dimensions:

\begin{table}[H]
\centering
\begin{tabular}{lll}
\toprule
\textbf{Dimension} & \textbf{Metric} & \textbf{Target} \\
\midrule
\multirow{3}{*}{Adaptation}
    & Learning Rate & $> 0.8$ in 100 interactions \\
    & Forgetting Rate & $< 0.1$ per week \\
    & Transfer Efficiency & $> 0.7$ cross-domain \\
\midrule
\multirow{3}{*}{Predictability}
    & Decision Variance & $< 0.15$ \\
    & Consistency Score & $> 0.9$ \\
    & Explanation Fidelity & $> 0.85$ \\
\midrule
\multirow{3}{*}{Structure}
    & Graph Sparsity & $0.3$--$0.7$ \\
    & Cluster Coherence & $> 0.8$ \\
    & Hierarchy Depth & $2$--$4$ levels \\
\bottomrule
\end{tabular}
\caption{PEMA-Bench metrics and targets}
\label{tab:pema-bench}
\end{table}

% ============================================================================
% PART III: EVALUATION
% ============================================================================

\part{Evaluation and Discussion}

\section{Evaluation}
\label{sec:evaluation}

\subsection{Experimental Design}

We evaluated \lumen across three enterprise deployments:

\begin{itemize}
    \item \textbf{Deployment A}: Financial services (500+ users, 18 months)
    \item \textbf{Deployment B}: Retail analytics (200+ users, 12 months)
    \item \textbf{Deployment C}: Manufacturing BI (150+ users, 6 months)
\end{itemize}

\subsection{Results}

\subsubsection{Task Success Rate}

\begin{table}[H]
\centering
\begin{tabular}{lccc}
\toprule
\textbf{System} & \textbf{Success Rate} & \textbf{Latency (p50)} & \textbf{Latency (p99)} \\
\midrule
Baseline (GPT-4 direct) & $71.4 \pm 3.2\%$ & $3.2 \pm 0.8$s & $12.1 \pm 2.3$s \\
Single-agent RAG & $82.1 \pm 2.8\%$ & $2.1 \pm 0.4$s & $8.4 \pm 1.6$s \\
\lumen (full) & $\mathbf{94.2 \pm 1.9\%}$ & $\mathbf{1.4 \pm 0.3}$s & $\mathbf{4.2 \pm 0.9}$s \\
\bottomrule
\end{tabular}
\caption{Comparative performance results (mean $\pm$ std over 3 deployments, $n=850$ users, $p < 0.01$ vs baseline)}
\label{tab:results}
\end{table}

\subsubsection{Ablation Study}

\begin{table}[H]
\centering
\begin{tabular}{lcc}
\toprule
\textbf{Configuration} & \textbf{Success Rate} & \textbf{$\Delta$ vs Full} \\
\midrule
Full \lumen & 94\% & --- \\
Without \pema & 89\% & -5\% \\
Without \sgrag & 85\% & -9\% \\
Without Two-Layer Routing & 91\% & -3\% \\
Without HallucinationGuard & 88\% & -6\% \\
\bottomrule
\end{tabular}
\caption{Ablation study results}
\label{tab:ablation}
\end{table}

\subsubsection{Failure Analysis}

We analyzed the 5.8\% failure cases ($n=49$ of 850 evaluated queries):

\begin{itemize}
    \item \textbf{Ambiguous queries} (41\%): User intent unclear, e.g., ``show me the numbers''
    \item \textbf{Schema mismatches} (31\%): Query references non-existent measures or dimensions
    \item \textbf{Complex multi-hop} (18\%): Queries requiring $>3$ agent handoffs
    \item \textbf{Token limits} (10\%): Context exceeded model limits
\end{itemize}

These failure modes inform future work on disambiguation prompts and chunking strategies.

\section{Discussion}
\label{sec:discussion}

\subsection{Trade-offs}

Key trade-offs identified:

\begin{itemize}
    \item \textbf{Latency vs Accuracy}: Deeper reasoning improves accuracy but increases latency
    \item \textbf{Specialization vs Generalization}: More agents enable specialization but increase coordination overhead
    \item \textbf{Plasticity vs Stability}: Higher learning rates enable faster adaptation but risk oscillation
\end{itemize}

\subsection{Limitations}

\begin{enumerate}
    \item \textbf{Domain Specificity}: Patterns optimized for BI may not transfer to other domains
    \item \textbf{Scale Bounds}: Not tested beyond 1000 concurrent users
    \item \textbf{Language Dependency}: Primarily evaluated on English and Spanish queries
\end{enumerate}

\subsection{Threats to Validity}

\textbf{Internal Validity}: Observed improvements may be influenced by dataset-specific optimizations. We mitigate this through ablation studies (Table~\ref{tab:ablation}) isolating each component's contribution.

\textbf{External Validity}: Results derive from three BI deployments in financial, retail, and manufacturing domains. Generalization to other domains (healthcare, legal) requires further validation.

\textbf{Construct Validity}: Task success relies on human evaluation. We employed two independent annotators per query with inter-rater reliability $\kappa = 0.84$.

\subsection{Ethical Considerations}

\begin{itemize}
    \item \textbf{Transparency}: All routing decisions are logged and explainable
    \item \textbf{Bias Monitoring}: Regular audits of agent behavior across user segments
    \item \textbf{Human Oversight}: Critical decisions require human confirmation
\end{itemize}

% ============================================================================
% CONCLUSIONS
% ============================================================================

\section{Conclusions}
\label{sec:conclusions}

This paper presented \lumen, a multi-agent agentic AI platform for Business Intelligence, along with three main contributions:

\begin{enumerate}
    \item \textbf{Functional Agency Framework}: Formal definitions establishing necessary conditions for production agentic systems

    \item \textbf{\pema Architecture}: Hebbian learning for adaptive agent routing with provable convergence bounds

    \item \textbf{Design Pattern Catalog}: 12 reusable patterns for enterprise AI systems
\end{enumerate}

Empirical evaluation across three enterprise deployments demonstrated:
\begin{itemize}
    \item 94\% task success rate (vs 71\% baseline)
    \item 40\% cost reduction through hybrid routing
    \item Hallucination reduction from 12\% to 0.3\%
\end{itemize}

\subsection{Future Work}

\begin{enumerate}
    \item Extension of \pema to cross-organization learning
    \item Integration of formal verification for safety-critical applications
    \item Development of domain-specific PEMA-Bench variants
\end{enumerate}

% ============================================================================
% APPENDICES
% ============================================================================

\appendix

\section{Formal Definitions}
\label{sec:formal-definitions}

\begin{definition}[Agent]
An agent $\alpha$ is a tuple $(\mathcal{S}, \mathcal{A}, \mathcal{T}, \pi, \rho)$ where:
\begin{itemize}
    \item $\mathcal{S}$: State space
    \item $\mathcal{A}$: Action space
    \item $\mathcal{T}$: Tool set
    \item $\pi: \mathcal{S} \rightarrow \mathcal{A}$: Policy function
    \item $\rho: \mathcal{S} \times \mathcal{A} \rightarrow \mathcal{S}$: Transition function
\end{itemize}
\end{definition}

\begin{definition}[Trust Graph]
The trust graph $G = (V, E, W)$ where:
\begin{itemize}
    \item $V = I \cup A$: Vertices are intents $I$ and agents $A$
    \item $E \subseteq I \times A$: Edges connect intents to capable agents
    \item $W: E \rightarrow [0,1]$: Weight function representing trust scores
\end{itemize}
\end{definition}

\section{Proofs}
\label{sec:proofs}

\begin{proof}[Proof of Theorem~\ref{thm:variance-bound}]

\textbf{Step 1: Single weight dynamics.}
Consider the weight update for a single intent-agent pair:
\[
w_{t+1} = w_t + \eta s_t - \lambda w_t = (1-\lambda)w_t + \eta s_t
\]
where $s_t$ is the success signal with $\mathbb{E}[s_t] = \mu$ and $\text{Var}[s_t] = \sigma^2$.

\textbf{Step 2: Stationary mean.}
At stationarity, $w^* = (1-\lambda)w^* + \eta\mu$, hence:
\[
\mathbb{E}[w_\infty] = \frac{\eta \mu}{\lambda}
\]

\textbf{Step 3: Variance recursion.}
The variance satisfies:
\[
\text{Var}[w_{t+1}] = (1-\lambda)^2 \text{Var}[w_t] + \eta^2 \sigma^2
\]
This is a geometric series with ratio $(1-\lambda)^2 < 1$ for $\lambda \in (0,1)$.

\textbf{Step 4: Equilibrium variance.}
At equilibrium, $\text{Var}[w_\infty] = (1-\lambda)^2 \text{Var}[w_\infty] + \eta^2\sigma^2$, yielding:
\[
\text{Var}[w_\infty] = \frac{\eta^2 \sigma^2}{1-(1-\lambda)^2} = \frac{\eta^2 \sigma^2}{2\lambda - \lambda^2} = \frac{\eta^2 \sigma^2}{\lambda(2-\lambda)}
\]

\textbf{Step 5: Multi-dimensional extension.}
For the weight vector $\mathbf{w} \in \mathbb{R}^{|I| \times |A|}$ with success signals having covariance matrix $\Sigma$ with maximum eigenvalue $\lambda_{\max}(\Sigma)$, the variance of each component is bounded by:
\[
\max_i \text{Var}[w_i] \leq \frac{\eta^2 \lambda_{\max}(\Sigma)}{\lambda(2-\lambda)}
\]

\textbf{Step 6: Policy variance bound.}
The routing policy $\pi_t = \text{softmax}(\mathbf{w}_t)$ maps weights to probabilities. By the delta method, $\text{Var}[\pi] \approx \|\nabla \pi\|^2 \text{Var}[w]$. Since softmax gradients satisfy $|\partial \pi_i / \partial w_j| \leq 0.25$ (maximum at $\pi_i = 0.5$), we have:
\[
\text{Var}[\pi_t] \leq 0.25 \cdot \text{Var}[w_t] \leq \frac{\eta^2 \sigma^2}{\lambda(2-\lambda)}
\]
where the factor 0.25 is absorbed into the inequality for simplicity.

\textbf{Stability:} The system is stable for any $\lambda \in (0,1)$ and $\eta > 0$ since $(1-\lambda)^2 < 1$ always holds. For fast convergence with low variance, we recommend $\eta \approx 0.01$ and $\lambda \approx 0.1$.
\end{proof}

\begin{remark}[Intuition]
The variance bound $\eta^2\sigma^2 / \lambda(2-\lambda)$ reveals a fundamental trade-off:
\begin{itemize}
    \item \textbf{Small $\eta$}: Low variance but slow adaptation
    \item \textbf{Large $\lambda$}: Faster forgetting reduces variance but limits memory
    \item \textbf{Sweet spot}: $\eta = 0.01$, $\lambda = 0.1$, $\sigma^2 = 0.25$ yields $\text{Var} \approx 0.00013$ (routing fluctuations $\approx 1.2\%$)
\end{itemize}
\end{remark}

\section{Reproducibility}
\label{sec:reproducibility}

\subsection{Hardware Configuration}

\begin{itemize}
    \item CPU: AMD EPYC 7V13 (32 cores)
    \item Memory: 128 GB DDR4
    \item GPU: NVIDIA A100 40GB (for embeddings)
    \item Storage: NVMe SSD
\end{itemize}

\subsection{Software Versions}

\begin{itemize}
    \item Python 3.11.4
    \item FastAPI 0.104.1
    \item Azure OpenAI API version 2024-02-15-preview
    \item Weaviate 1.23.0
    \item Redis 7.2
\end{itemize}

\subsection{Availability}

This whitepaper, diagrams, and LaTeX source available at:
\url{https://github.com/innerbi/lumen-whitepaper}

% ============================================================================
% REFERENCES
% ============================================================================

\bibliographystyle{plainnat}
\bibliography{references}

\end{document}
